\chapter{Running in the kernel}
\label{ch:kernel}

Everything so far has run in userspace, where a mistake is a segfault under
\code{gdb} and thirty seconds. This chapter is the first time \code{core/} runs
inside the kernel, where the same mistake is a panic and a reboot.

The claim being tested is the one the seam has been making since phase~1: that
if \code{core/} is honest about its dependencies, arriving in the kernel adds a
second environment implementation and changes nothing above it. That claim
turns out to hold, and the evidence arrived before a line of VFS code was
written --- see \S\ref{sec:kernel-buildproof}.

What phase~5 does \emph{not} produce is a working mount. The on-disk format has
no FS tree, so there is no root directory to build an inode from and nothing to
hand \fnname{d\_make\_root}. Mounting therefore ends in a refusal --- but a
refusal that has already read the superblock, validated it six ways, and walked
into the root tree through the new environment. Every step before the last one
is permanent.

\section{Two builds of one directory}
\label{sec:kernel-buildproof}

\code{core/} is compiled twice from the same sources: by \code{core/Makefile}
into \code{libbitter.a} for the tools and tests, and by Kbuild with the
kernel's own flags into \code{bitterfs.ko}. Nothing in it is conditional on
which.

Kbuild builds each object into a directory beneath \code{M=}, and a line
reading \code{bitterfs-y += ../core/btree.o} asks it to write outside the tree
it was told to manage. The resolution was to move \code{M=} \emph{up} to the
repository root rather than reach down from \code{kernel/}: every object is
then genuinely beneath it, and \code{-I\$(src)} resolves to the same directory
\code{core/Makefile} passes as \code{-I\$(TOP)}, so an include means the
identical thing in both builds.

The first such build is the proof. All eight \code{core/} files compiled under
kernel flags with no errors and no warnings, and the module's complete list of
undefined symbols was:

\begin{quote}
\code{\_printk}, plus compiler-inserted KASAN, stack-protector and retpoline
entries.
\end{quote}

\code{core/} requires nothing from the kernel --- not even \fnname{memcpy},
since \code{core/bitter\_string.h} supplies its own. Four phases of
\code{make freestanding} bought exactly what that target claimed it would.

The cost was discovered immediately afterwards, and is recorded in
\code{docs/LOG.md}: both builds wanted to write \code{core/btree.o}, and the
kernel's object is \code{-mcmodel=kernel}, non-PIE and KASAN-instrumented.
Archiving it into \code{libbitter.a} fails the userspace link with a message
that names neither the kernel nor the cause. The userspace build now writes to
\code{core/build/}.

\section{Stopping}

\code{core/bitter\_assert.h} declares \fnname{bitter\_assert\_fail} and leaves
each consumer to define it: \code{user/util.c} prints and calls
\fnname{abort}, and \code{kernel/assert.c} is the kernel's answer. A link-time
seam rather than an entry in \code{bitter\_env\_ops}, because there is nothing
per-filesystem to configure.

The header specifies \code{\_Noreturn}, and that is a promise rather than a
mechanism: the compiler deletes the fall-through path at every call site. A
kernel implementation that warned and continued would therefore not merely
behave oddly, it would execute whatever bytes followed the call.

\code{BUG()} satisfies the promise honestly. On x86 it is three lines --- an
\code{ud2} planted inline, and \fnname{\_\_builtin\_unreachable} to tell the
compiler so. The trap handler prints, dumps registers and a stack trace, and
kills the \emph{task}.

\paragraph{Why the function prints first.} \code{CONFIG\_DEBUG\_BUGVERBOSE}
prints a file and line of its own, but they come from \code{\_\_bug\_table} and
name the site of the \code{BUG()} macro --- which is \code{kernel/assert.c},
every time, whichever invariant actually failed. The \field{expr},
\field{file} and \field{line} arguments are the only carriers of the real
location, and nothing prints them for you.

\decide{Task-fatal was chosen over \fnname{panic}, so the guest survives long
enough to read \code{dmesg}. The cost is that nothing unwinds: the
filesystem-wide mutex is never released and the mount hangs, with
\code{CONFIG\_DETECT\_HUNG\_TASK} burying the message under a splat every 120
seconds. Reconsider when an assertion fires somewhere reboot-hostile.}

\section{The second environment}

\code{kernel/env\_kernel.c} supplies the six operations of
\code{bitter\_env\_ops} over the buffer cache, as \code{user/env\_user.c}
supplies them over \fnname{pread} and \fnname{malloc}.

Read-only reaches two of them. \fnname{alloc\_block\_buf},
\fnname{dirty\_block}, \fnname{flush\_device} and \fnname{writeback\_all} are
all write-side, and each is a \code{BITTER\_ASSERT(0)}: \code{core/} arriving
there is a bug in the module rather than anything the disk did, which is the
test \code{bitter\_assert.h} states. Each still returns afterwards, because
\code{BITTER\_ASSERT} compiles to nothing without \code{BITTER\_DEBUG} and a
non-void function must not fall off its end.

\fnname{bitterfs\_read\_block} is where the work is. It fetches through
\fnname{sb\_bread}, which is a cache lookup first and a synchronous read only
on a miss, and then applies the same three self-identifying checks
\code{env\_user} applies, in the same order and for the same reason: checksum
says the bytes are intact, \field{bytenr} says it is the block that was asked
for, \field{fsid} says it belongs to this filesystem rather than to a previous
\fnname{mkfs} at the same address. Each presupposes the one before it.

Three differences from the userspace version are worth stating.

\paragraph{No copy.} \field{b\_data} points \emph{into} the page cache. The
buffer head is a descriptor; the bytes live in a cached page, and copying them
would decouple what \code{core/} modifies from what is eventually written.

\paragraph{Allocation comes last.} \code{env\_user} must allocate first,
because it reads into its own storage. \fnname{sb\_bread} hands over the data,
so a block that fails a check costs no allocation at all --- and there is one
fewer failure path to unwind.

\paragraph{No \field{have\_fsid} gate.} \code{env\_user} gates its fsid check
because \fnname{mkfs} and the unit fixtures use the seam with no superblock at
all. The kernel has no such caller: every read happens after
\fnname{fill\_super} has read the superblock, because \fnname{fill\_super} is
what installs the vtable. The invariant is enforced by what is reachable rather
than by a flag --- but it is then \fnname{fill\_super}'s to keep, and it is
written down on the struct in \code{kernel/bitterfs.h}. An unarmed fsid is
sixteen zero bytes, which would silently match a zeroed block rather than fail.

\decide{\code{core/} has never run against a caching environment. Two
\fnname{read\_block} calls for one address alias in the kernel and do not in
userspace, and every test uses the latter. Harmless while nothing writes. See
\code{docs/LOG.md}.}

\section{Mounting}

\code{kernel/super.c} uses the \code{fs\_context} mount API --- an
\fnname{init\_fs\_context} that installs an operations table, a
\fnname{get\_tree} that calls \fnname{get\_tree\_bdev}, and a
\fnname{reconfigure} that reasserts read-only. The legacy \fnname{mount\_bdev}
form would have been three lines instead of twelve; \code{fs\_context} was
chosen because mount options arrive through a typed \fnname{parse\_param}
rather than an unparsed blob, and because \fnname{errorfc} reports failures to
whoever ran \code{mount}. That last point decided it: for the whole of phase~5,
\fnname{fill\_super}'s job \emph{is} to refuse and say why.

\subsection{Ordering}

\fnname{fill\_super} has five constraints, and each exists because something
downstream reads what the previous step wrote.

\begin{enumerate}
\item \fnname{sb\_set\_blocksize} before any read, because \fnname{sb\_bread}
      counts in \field{sb->s\_blocksize}. It returns \code{0} on failure ---
      the opposite convention to \fnname{set\_blocksize} beneath it, so the
      check is truthiness and not \code{< 0}.
\item \field{s\_fs\_info} assigned immediately after the allocation, which
      hands ownership of the mount struct to \fnname{kill\_sb}. From that line
      on, no error path frees it; doing so would be a double free.
\item The superblock read \emph{outside} the seam. \fnname{read\_block} cannot
      read it: the checksum covers \bytes{512} rather than the whole block, the
      layout is \code{bitter\_super} rather than \code{bitter\_header} so the
      field offsets differ, and the fsid it would check against is the value
      this read supplies.
\item Fields copied out before the \fnname{brelse}, because the pointer aims
      into the page cache.
\item \fnname{bitterfs\_env\_init} only after the fsid is stored --- see the
      previous section.
\end{enumerate}

\paragraph{Cleanup is asymmetric.} \fnname{put\_super} does not run when
\fnname{fill\_super} fails, because \fnname{generic\_shutdown\_super} guards it
with \code{if (sb->s\_root)} and \field{s\_root} is never set on a failed
mount. \fnname{kill\_sb} \emph{does} run. That is why the freeing lives there
and not in a \code{super\_operations} method: it is the one that runs both ways.

\subsection{Six checks and one label}

Magic first, and silently: this is the ``not mine'' answer, and \code{mount}
may be probing several filesystem types in turn. \code{-EINVAL} with no text is
how a filesystem declines a device politely.

Every check after it means ``this \emph{is} a bitterfs and something is wrong
with it,'' which is when the caller needs to be told what. Checksum,
\field{bytenr}, \field{block\_size} and \field{root\_level} report
\code{-EUCLEAN}; \field{csum\_type} reports \code{-EOPNOTSUPP}, because the
field is valid and names an algorithm this build does not implement rather than
describing damage. That is the same category the feature flags will fall into,
and they will want checking in the same place.

All six exit through one \code{out\_brelse} label. \code{core/}'s
\fnname{bitter\_path\_release} makes the argument for why: the paths that owe a
release are precisely the ones least likely to be exercised.

\section{What the first mount proved}

\begin{quote}
\code{bitterfs: no FS tree in this image (-2)}
\end{quote}

\code{-2} is \code{-BITTER\_ENOENT}, returned by \fnname{bitter\_find\_root}.
Reaching it required setting the block size, reading the superblock at block
16, passing all six checks, arming the seam, and then calling into
\code{core/} --- which went through \fnname{bitterfs\_read\_block}, fetched the
root tree block, verified its checksum, \field{bytenr} and \field{fsid},
searched the leaf, and found no \itemtype{ROOT\_ITEM} for
\code{BITTER\_FS\_TREE\_OBJECTID}. \fnname{put\_block} then released it, and
\code{rmmod} was clean with KASAN watching, so nothing leaked.

The error is reported as \code{-EUCLEAN} rather than passed through.
\code{BITTER\_ENOENT} is 2, which is \code{ENOENT}, and a mount failing with
``no such file or directory'' reads as though the \emph{device} were missing.

The automated check in \code{qemu/guest.sh} asserts that exact message and that
exact number, so any earlier failure --- a bad checksum, a foreign fsid, an I/O
error --- fails the test rather than passing quietly.

\decide{The mount ends here until the on-disk format grows an FS tree: inode
items, directory entries, and a root directory for \fnname{mkfs} to create.
That is userspace work, and it is what phase~5 has left.}
